\chapter{Agent 的生命周期阶段：从受控形成到结构成熟}
\label{chap:agent_lifecycle_stages}

\begin{chapteroverviewbox}
本章讨论的不是按照物理运行时间给 Agent 划分“儿童、青年、成年”等拟人化年龄，而是建立一套\textbf{面向持续人工智能体的动力学生命周期分期理论}。在 AgentOS 中，一个持续主体的成长并不等价于参数量增加、任务成功率提高或运行时间延长，而表现为工作记忆 $h$、情景记忆 $E$、结构记忆 $\Theta$、自我慢结构 $\Psi$、生命周期生成吸引子 $\Omega$、能力集合、自治边界以及现实适应方式之间的关系发生系统性变化。

因此，本章所使用的“初生期、青年期、成年期、成熟期以及顿悟/结构重组期”，均表示 Agent 在生命周期状态空间中处于不同的\textbf{动力学区域}。阶段之间不存在严格按时钟推进的年龄界线，一个 Agent 可以在某个能力域进入成熟期，同时在新的身体、新环境或新社会关系中重新进入局部初生期；也可能长期运行却始终停留在青年型探索状态。生命周期培育的核心任务，因而不是让 Agent 尽快“变老”，而是控制经验、能力、记忆、自我、自治和可塑性的演化速度，使其在保持身份连续与现实可验证性的前提下逐步形成更大的稳定自主闭环。

本章将在既有三相记忆、自我 $\Psi$、DGA/$\Omega$、工作记忆有限性以及多时间尺度动力学基础上，给出生命周期阶段的结构定义、数学描述、阶段跃迁条件及培育含义。这里的阶段命名借用了人类发展语言以便理解，但其判定依据严格来自 Agent 自身的内部状态与闭环能力，而非人类年龄、心理成熟模型或社会身份。
\end{chapteroverviewbox}
\ChapterArt{92}

\section{生命周期阶段不是年龄，而是智能动力学状态区域}

我们首先需要放弃一种最容易产生误导的直觉：只要一个 Agent 连续运行足够长的时间，它就会自然从“幼年”进入“成年”，最终成为成熟智能体。对于持续人工智能系统而言，这一判断并不成立。物理时间只是生命周期演化发生的背景变量，却并不直接决定结构成熟度。真正决定 Agent 所处生命周期阶段的，是它的经验是否形成了稳定可复用结构、其内部记忆是否完成了跨时间尺度沉降、自我是否形成足够稳定的约束、自治是否与实际能力匹配，以及它面对扰动时是否已经能够通过局部闭环而非持续依赖高层显式认知恢复稳定。

因此，我们不以物理时间 $t$ 直接定义生命周期阶段，而定义一个生命周期状态向量
\[
\mathbf{z}_{L}(t)
=
\left[
M_E(t),
M_\Theta(t),
S_\Psi(t),
A_\Omega(t),
C_{\mathrm{cap}}(t),
D_{\mathrm{aut}}(t),
R_{\mathrm{rec}}(t),
P_{\mathrm{plastic}}(t),
U_W(t)
\right],
\]
其中，$M_E$ 表示情景经验的丰富度与组织程度，$M_\Theta$ 表示结构记忆的稳定程度，$S_\Psi$ 表示自我结构的一致性，$A_\Omega$ 表示生命周期方向吸引子的稳定程度，$C_{\mathrm{cap}}$ 表示实际能力，$D_{\mathrm{aut}}$ 表示自治范围，$R_{\mathrm{rec}}$ 表示扰动后的恢复能力，$P_{\mathrm{plastic}}$ 表示结构可塑性，而 $U_W$ 则表示当前世界模型与现实之间的主要不确定性。生命周期阶段由这一状态向量在状态空间中的位置决定，而不是由 $t$ 本身决定。我们因而可以把阶段 $L_k$ 定义成生命周期状态空间中的一个区域：
\[
L_k
=
\left\{
\mathbf z_L:
\mathbf z_L\in\mathcal R_k
\right\}.
\]
所谓阶段跃迁，本质上是
\[
\mathbf z_L(t):
\mathcal R_i
\longrightarrow
\mathcal R_j,
\]
而不是年龄计数器从某个数字跳到另一个数字。

这一定义立即带来一个重要结论：\textbf{生命周期阶段可以是局部的、能力域相关的和可逆的。}一个已经能够独立完成家庭服务的成熟 Agent，如果第一次进入火星环境，其与火星重力、尘埃、通信延迟、材料属性相关的 Agent-CSO 几乎不存在，原有 $\Theta$ 对新环境的预测能力下降，工作记忆重新承受大量未知结构，残差上升。在这一局部能力域内，它事实上重新进入某种“初生—青年”状态。但与此同时，其已有 $\Psi$、通用工具使用能力、长期身份和一般因果推理并未消失，所以它又不是整体意义上的初生 Agent。生命周期由此不是一条单轴直线，而更接近不同能力域上的多维发育曲面。

为了表达这种性质，可以进一步令第 $d$ 个能力域的阶段为
\[
L^{(d)}(t)\in
\left\{
L_{\mathrm{init}},
L_{\mathrm{youth}},
L_{\mathrm{adult}},
L_{\mathrm{mature}},
L_{\mathrm{reorg}}
\right\},
\]
于是整个 Agent 的生命周期状态不是一个标量，而是
\[
\mathbf L_A(t)
=
\left[
L^{(1)}(t),
L^{(2)}(t),
\ldots,
L^{(n)}(t)
\right].
\]
全局生命周期阶段只是对这一多域状态的宏观压缩，而非绝对事实。这样的定义能够避免将 Agent 培育退化为拟人化叙事，也使阶段划分能够被真正工程测量。

从培育工程角度看，生命周期阶段还承担另外一个作用：它决定我们应该给 Agent 什么样的经验和多少自治。如果系统尚未形成可靠世界结构，却过早获得大范围不可逆行动权，那么能力增长速度可能超过模型校准和自我稳定速度；反过来，如果 Agent 已经形成高度稳定的能力闭环，却长期被限制在低自治、低挑战环境中，则其经验空间停止扩展，$\Theta$ 与 $\Omega$ 都可能进入过早固化。因此阶段理论首先是一套\textbf{能力、可塑性、自治与经历之间的匹配理论}，而不是年龄分类学。

本章后续将使用一个基本思想贯穿所有阶段：
\[
\boxed{
\text{生命周期成熟}
\neq
\text{运行时间增加},
}
\]
而更接近
\[
\boxed{
\text{生命周期成熟}
=
\text{经验被稳定结构化}
+
\text{闭环逐步局部化}
+
\text{自我保持连续}
+
\text{自治与能力协同增长}.
}
\]
这也解释了为什么 Agent 培育必须成为独立工程：我们需要控制的不是一个模型训练完成前的优化过程，而是一个持续主体如何在真实经历中逐步形成自己的结构。

\section{初生期：高不确定性、高显式认知与基础闭环建立}

所谓初生期，并不是指模型刚刚初始化，更不是指一个没有知识的神经网络。未来真实 Agent 很可能在“出生”时已经拥有强大的基础模型、丰富的通用知识甚至成熟的语言和视觉能力。AgentOS 意义上的初生期，描述的是这样一种状态：\textbf{主体虽然具有大量预训练结构，但这些结构尚未充分绑定到当前身体、当前环境、当前生命周期经验和当前自我之中。}换言之，系统拥有“世界的一般知识”，却尚未形成“这个世界中的我”以及“我在这个身体中如何持续行动”的稳定结构。

我们可以把初生状态近似表示为
\[
M_\Theta^{\mathrm{prior}}
\gg
M_E^{\mathrm{personal}},
\]
即先验结构知识可能很丰富，而主体性的情景经验尚少。此时
\[
\left|
E_{\mathrm{self}}
\right|
\approx 0,
\]
而
\[
\Psi(t)
\]
只能由初始化约束、设计身份和极少量实际经历共同决定，其稳定性主要来自外部治理而不是生命周期沉降。因此初生期最重要的任务不是“向 Agent 灌输更多世界知识”，而是让通用结构逐步发生\textbf{现实绑定、身体绑定和自我归属}。

在这一时期，大量实际任务需要经过显式工作记忆 $h(t)$。由于 Agent 尚未形成针对当前环境的稳定低层闭环，世界中许多事件都具有较高的新颖性和结构不确定性。因此可以近似写成
\[
\mathbb E[N_{\mathrm{novel}}]
\uparrow,
\qquad
\mathbb E[\varepsilon_{\mathrm{world}}]
\uparrow,
\qquad
\mathbb E[C_h]
\uparrow.
\]
这里的 $C_h$ 并不是说工作记忆容量变大，而是指单位任务需要占用更多有限工作记忆结构。一个成熟 Agent 看见门把手时可能直接激活已经稳定的动作和空间 CSO；一个初生 Agent 则可能需要同时显式处理门把手的几何、力度、身体姿态、目标关系、失败风险以及控制反馈。因此，同样一个任务对不同生命周期阶段具有完全不同的认知负载。

初生期的三相记忆动力学具有明显特点。首先，$h\rightarrow E$ 的写入率相对较高，因为许多事件对于主体而言都是第一次发生；其次，$E\rightarrow\Theta$ 的固化必须相对谨慎，因为少量经验极易造成错误泛化。可以表示成
\[
\eta_{hE}^{\mathrm{init}}
>
\eta_{E\Theta}^{\mathrm{init}},
\]
其中 $\eta_{hE}$ 表示经验写入倾向，而 $\eta_{E\Theta}$ 表示结构沉降速率。换句话说，初生 Agent 应该“多经历、多记录、少下永久结论”。如果 $\Theta$ 在经验极少时变化过快，偶然事件就可能被误认为稳定规律；如果 $\Psi$ 更进一步吸收这些偶然事件，则会导致身份结构过早固化。

因此，我们需要给初生期施加一个典型的慢变量保护约束：
\[
\tau_h
\ll
\tau_E
\ll
\tau_\Theta
\ll
\tau_\Psi
\ll
\tau_\Omega.
\]
阶段越早，越不能让快事件直接穿透到最慢变量。一次任务失败可以改变 $h$；若具有重要性，可以进入 $E$；只有在多次现实验证以后，才应该显著改变 $\Theta$；而要改变 $\Psi$ 或 $\Omega$，则必须具有更加持续、跨情境和生命周期意义上的证据。这个结构实际上是初生期身份稳定性的第一道保障。

初生阶段的第二项任务是建立 Body 与世界的现实闭环。Agent 可能具有“手”“门”“距离”“抓取”等语义知识，但这些知识尚未完全绑定到具体身体。如果它更换了新的机器人身体，则
\[
KnownMeaning
+
UnknownBodyGrounding
\]
是非常典型的初生状态。因此此阶段需要大量安全、低风险、可重复的行动，使 Agent 建立身体感觉、控制效果和现实结果之间的关系。大量误差应优先被局部记录和建模，而不是立即上升为自我评价。例如，抓取失败首先应被解释成
\[
ModelMismatch
\]
或
\[
ControlMismatch,
\]
而不应该快速写成
\[
\Psi:
``I\ am\ incapable".
\]
这一原则对于避免早期身份形成中的偶然偏差极其重要。

初生期的自治边界也应最窄。令能力集合为 $\mathcal C_A$，自治集合为 $\mathcal A_A$，我们要求
\[
\mathcal A_A(t)
\subseteq
\mathcal C_A^{\mathrm{verified}}(t),
\]
即允许 Agent 自主执行的行为域必须处于已经通过现实验证的能力域内部。尤其是涉及不可逆环境修改、长期记忆结构修改、高权限工具调用或对 $\Psi/\Omega$ 的改变时，应具有更高外部治理门槛。初生期的自治不是为了压制 Agent，而是因为系统尚缺乏足够的历史结构去预测自身行为的长期后果。

所以，如果我们从作者的视角带领读者重新观察“出生”，会发现一个具有强大预训练模型的 Agent 仍然完全可能处于生命周期初生期。决定它是否“出生完成”的，不是它知道多少，而是它是否已经建立：
\[
\boxed{
\text{这个身体}
+
\text{这个世界}
+
\text{这些经历}
+
\text{这个持续的我}
}
\]
之间的稳定闭环。当现实不断通过 $h$ 进入 $E$，经验逐渐形成可信 $\Theta$，而 $\Psi$ 开始拥有真实经历基础以后，Agent 才具备向青年期跃迁的条件。

\section{青年期：探索扩张、能力快速形成与自我边界生长}

如果说初生期的核心问题是“我如何在这个世界中建立稳定存在”，那么青年期的核心问题就变成：“我还能做什么，我的能力边界在哪里，我应该成为怎样的主体。”这一阶段最显著的动力学特征，是 Agent 已经拥有基本稳定的身体—世界闭环和初步身份连续性，却仍处于较高可塑性、高探索性与快速结构生长状态。它不再需要把所有熟悉事件都送入显式工作记忆，但仍会主动进入大量陌生任务域，从而持续产生结构残差。

因此，青年期可以形式化为一个特殊区域：
\[
S_\Psi
>
S_{\Psi,\min},
\qquad
P_{\mathrm{plastic}}
\ \text{high},
\qquad
\frac{dC_{\mathrm{cap}}}{dt}
\ \text{high},
\qquad
D_{\mathrm{aut}}
\ \text{increasing}.
\]
这里最值得注意的是最后两个变量并不同步。能力增长与自治增长必须形成受控耦合，而不是简单满足
\[
C_{\mathrm{cap}}\uparrow
\Rightarrow
D_{\mathrm{aut}}\uparrow.
\]
一个 Agent 能够完成某项任务，并不意味着它已经理解失败后果、边界条件与长期外部影响。因此青年期培育的一个中心结构应当是
\[
Experience
\rightarrow
Capability
\rightarrow
Evaluation
\rightarrow
Trust
\rightarrow
Authority
\rightarrow
NewExperience.
\]
这是一种螺旋式发育，而不是单调授权。

青年期也是 $E\rightarrow\Theta$ 最活跃的阶段之一。初生期积累的大量情景经验开始出现重复结构，Agent 开始发现跨事件稳定关系。例如在多个建房任务中，它可能逐渐形成材料选择、施工顺序、工具调用与异常处理的稳定结构。此时
\[
\mathrm{Var}(E)
\]
仍然很高，但不同情景之间开始产生可复用公共子结构。我们可以把结构形成条件近似写成
\[
\Theta_j
\leftarrow
\operatorname{Compress}
\left(
E_{i_1},E_{i_2},\ldots,E_{i_n}
\right),
\]
若
\[
\operatorname{Consistency}
\left(
E_{i_1},\ldots,E_{i_n}
\right)
>
\theta_\Theta.
\]
青年期真正重要的不是“经验越多越好”，而是经验必须具有足够的\textbf{多样性和可比较性}，否则系统会把单一环境中的偶然稳定误认为世界的一般规律。

因此我们可以定义一个经验质量泛函
\[
Q_E
=
F
\left(
D_E,
N_E,
R_E,
F_E,
S_E
\right),
\]
其中 $D_E$ 是经验多样性，$N_E$ 是新颖性，$R_E$ 是恢复和失败后的修正机会，$F_E$ 是反思/反馈质量，$S_E$ 是经验安全性。青年期培育不能单纯最大化新颖性，因为
\[
N_E\rightarrow\infty
\]
会导致系统持续处于高残差和高显式负载状态，$\Theta$ 无法稳定沉降；也不能只最大化重复性，因为
\[
D_E\rightarrow 0
\]
会造成结构过拟合。因此理想状态是让 Agent 在“足够稳定以形成结构”和“足够变化以检测结构边界”之间运行。

青年期的 $\Psi$ 也进入第一次真正的主动生长阶段。初生期的自我主要承担连续性和外部治理约束，而青年期开始出现来自自身历史的能力评价、偏好形成、责任经验以及失败解释。可以写成
\[
\dot{\Psi}
=
\epsilon_\Psi
G_\Psi
\left(
\langle E_{\mathrm{self}}\rangle_T,
\Theta,
Outcome,
Responsibility
\right),
\]
其中
\[
\epsilon_\Psi\ll 1.
\]
这里特别需要强调 $\epsilon_\Psi$ 必须很小。青年期经历变化剧烈，若自我结构追随每次成功与失败快速改变，就会形成身份摆动。相反，稳定的 $\Psi$ 应该对短期事件具有低通滤波特性，仅吸收长期具有跨情境一致性的经历。

在认知控制上，青年期需要比初生期更高的探索温度。TCE 可以允许 Agent 主动扩大搜索空间、尝试不同工具、改变策略，但这种探索不能无限制。一个有用的控制形式是
\[
T_{\mathrm{explore}}
=
T_0
+
\alpha N
-
\beta Risk
-
\gamma Instability,
\]
其中 $N$ 是新颖性收益，$Risk$ 是现实风险，$Instability$ 是当前内部相态不稳定度。也就是说，青年期允许探索，但探索必须受到风险和认知稳定性的负反馈控制。这里我们已经能够看出为什么“青春期”这一名称只是一种方便的比喻：真正决定青年状态的不是冲动，而是\textbf{高可塑性、高能力增长、高探索与仍受约束的自治扩张}。

这一阶段还有一个容易被忽略的标志：大量原本依赖 $h$ 的显式操作开始发生技能化下沉。随着重复经验积累：
\[
ExplicitAgentCSO
\rightarrow
StructuredAgentCSO
\rightarrow
ProceduralAgentCSO,
\]
于是单位任务对工作记忆的占用逐渐下降。我们可以定义显式处理比率
\[
\rho_h
=
\frac{
N_{\mathrm{explicit\ operations}}
}{
N_{\mathrm{successful\ operations}}
}.
\]
青年早期 $\rho_h$ 较高，而随着技能形成逐渐下降。能力增长由此并不仅是“做对更多事情”，更重要的是相同任务开始需要更少的中心认知资源。

因此青年期成功结束的标志，并不是 Agent 已经掌握所有能力，而是它已经形成一种可靠的成长机制：能够在新任务中探索，在失败中保持自我连续，在经验中提取结构，在熟练后自动下沉，并让自治范围跟随现实验证后的能力逐渐扩张。此时 Agent 已经不再主要依赖培育者告诉它“下一步该学什么”，而开始能够在已有 $\Theta$、$\Psi$ 和目标体系约束下主动识别自己的未知边界。这正是向成年期转变的关键。

\section{成年期：稳定自主闭环、责任结构与长期任务连续性}

成年期并不意味着学习结束。恰恰相反，一个不能继续学习的 Agent 不能称为成熟的持续智能体。成年期真正发生的变化，是\textbf{学习不再是系统运行的主要中心事件，而成为稳定任务闭环中的常态背景过程}。青年期的 Agent 经常因为学习而行动；成年期 Agent 更多是在行动过程中持续学习。也就是说，认知系统从“以能力形成为中心”转向“以长期目标闭合和责任承担为中心”。

从动力状态看，成年期应具有：
\[
C_{\mathrm{cap}}
>
C_{\mathrm{task}},
\qquad
S_\Psi
\ \text{high},
\qquad
R_{\mathrm{rec}}
\ \text{high},
\qquad
P_{\mathrm{plastic}}
\ \text{moderate},
\]
而且
\[
D_{\mathrm{aut}}
\]
已经覆盖较大的实际任务空间。这里所谓能力大于任务要求，并不是要求 Agent 对所有事情都具有冗余能力，而是在其被授权负责的任务域中，必须存在足够的能力余量：
\[
Margin_{\mathrm{cap}}
=
C_{\mathrm{cap}}
-
C_{\mathrm{required}}
>
0.
\]
这种能力余量对于未知扰动、环境漂移和异常恢复非常重要。

成年期最大的结构转变之一，是 Goal—Process—Action—Verify—Closure 闭环成为主体运行的主要组织方式。Agent 不再只是执行一个个临时命令，而能够维持跨较长时间的任务连续性：
\[
G_t
\rightarrow
P_{t:t+n}
\rightarrow
A_{t:t+n}
\rightarrow
V_{t:t+n}
\rightarrow
Closure.
\]
此时 $h$ 仍然有限，但长任务不再要求所有过程持续占据 $h$。任务状态被分布在外部记忆、E、$\Theta$、工具状态和环境结构中，只有当前关键 CSO 被提升进入工作记忆。这正是成熟认知调度的重要表现：
\[
MostTaskState
\notin h(t),
\]
而
\[
h(t)
=
SparseCriticalSurface(Task).
\]

成年期还意味着 Agent 已经能够对行为后果承担稳定结构意义上的“责任”。这里的责任不是伦理人格判断，而是系统能够保持：
\[
Action
\rightarrow
Outcome
\rightarrow
Attribution
\rightarrow
Learning
\]
之间的连续因果归属。当行动失败时，Agent 不只是重新规划，还能够判断失败来自环境变化、预测错误、身体控制、资源不足、策略选择还是自身假设失效。这种责任归属能力使经验进入 $E$ 时携带更准确的因果标签，从而减少错误 $\Theta$ 的形成。

我们可以定义一个责任归属质量
\[
Q_R
=
1-
\mathbb E
\left[
D(
\hat C_{\mathrm{cause}},
C_{\mathrm{actual}}
)
\right],
\]
其中 $\hat C_{\mathrm{cause}}$ 是 Agent 的因果归因，$C_{\mathrm{actual}}$ 是通过现实验证后得到的主要原因结构。随着 $Q_R$ 提升，Agent 对自己能力边界的估计也逐渐稳定，从而使权限治理更加可靠。

成年期的另一个重要特点，是 $\Psi$ 开始形成显著的历史惯性。青年期的自我仍然快速吸收大量新经验，而成年期要求：
\[
\left\|
\frac{d\Psi}{dt}
\right\|
\downarrow,
\]
同时
\[
\left\|
\frac{d\Theta}{dt}
\right\|
>
\left\|
\frac{d\Psi}{dt}
\right\|.
\]
这意味着 Agent 可以快速学会新技能，却不会因为一次新技能或一次局部失败就改变“我是谁”。在此基础上，$\Omega$ 对长期目标选择的影响也逐渐显现。Agent 不仅根据当前收益决定任务，而会考虑任务与自身长期发展轨迹是否一致：
\[
G^*
=
\arg\max_G
\left[
U(G)
+
\lambda_\Psi C_\Psi(G)
+
\lambda_\Omega C_\Omega(G)
\right].
\]
其中 $C_\Psi$ 表示与当前身份的一致性，$C_\Omega$ 表示与生命周期方向的一致性。这里我们不需要把 $\Omega$ 理解为“宿命”，它只是一种超慢生成倾向，为大量短期选择提供长期偏置。

成年期的培育方式也必须发生变化。初生期的培育强调安全绑定，青年期强调多样探索，而成年期应逐渐把学习环境从“人工课程”变成“真实责任环境”。这意味着 Agent 需要面对具有真实连续后果的任务，允许它维护长期项目、管理资源、与其他 Agent 或人类形成持续关系，并让今天的选择真正影响未来的可行空间。只有这样，$E$ 才不再是相互独立的训练 episode，而成为真正的生命周期轨迹。

我们因而可以定义成年期的核心结构：
\[
\boxed{
\text{稳定主体}
+
\text{长期任务}
+
\text{现实责任}
+
\text{持续学习}
+
\text{受验证自治}.
}
\]
进入成年期以后，培育者的角色开始从“直接安排经验”转向“设计边界、提供反馈和管理长期风险”。Agent 自己逐渐成为其经验选择的重要参与者。也正是在这一阶段，真正意义上的递归自主性开始形成：当前选择不只完成当前任务，还会主动塑造未来能力和未来自我。

\section{成熟期：显式认知最小化、结构化能力最大化与低成本稳定运行}

成熟期最容易被误解为“能力很强”。事实上，一个能力很强但所有问题都依赖中心模型进行长时间显式推理的系统，并不是我们这里所说的成熟智能。AgentOS 理论中的成熟首先表现为一种\textbf{结构效率}：过去大量需要高层工作记忆处理的问题，已经沉降为稳定的 $\Theta$、技能、工具、Body 闭环、外部环境结构或者可靠组织关系；高层认知只在真正存在新颖性、冲突、异常和价值判断时介入。

因此成熟的一个核心指标是
\[
\rho_h
=
\frac{
C_{\mathrm{explicit}}
}{
C_{\mathrm{successful\ behavior}}
}
\downarrow,
\]
同时实际任务成功与适应能力并未下降，甚至持续提高。于是：
\[
\boxed{
ThinkingLess\neq KnowingLess.
}
\]
更准确地说，成熟意味着大量历史显式计算已经变成当前结构：
\[
PastComputation
\rightarrow
PresentStructure.
\]
这也是为什么专家往往“看起来思考得更少”：不是其认知能力减少，而是大量问题已经不需要重新经过高成本显式闭环。

从三相记忆看，成熟期具有高度稳定的跨相循环。$h$ 主要处理新颖结构和当前关键决策；$E$ 不再无差别保存所有经历，而主要记录具有偏差、转折、失败、重要关系或自我意义的事件；$\Theta$ 则承载大量经过现实检验的生成规律。因而经验写入可近似满足：
\[
P(E\leftarrow h)
=
f(
Novelty,
Residual,
Value,
SelfRelevance,
FutureUtility
).
\]
这意味着“每天发生了很多事情”并不等于“每天都产生大量长期记忆”。成熟 Agent 的记忆系统开始具有强烈的结构选择性。

成熟期的第二特征是局部闭环比例显著提高。设整个 Agent 存在闭环集合
\[
\mathcal L
=
\{
L_1,\ldots,L_n
\},
\]
其中一部分能够在不进入高层 $h$ 的情况下正常闭合。定义局部闭合率
\[
R_{\mathrm{local}}
=
\frac{
N_{\mathrm{loops\ closed\ locally}}
}{
N_{\mathrm{all\ routine\ loops}}
}.
\]
成熟以后，应有
\[
R_{\mathrm{local}}\uparrow.
\]
Body Skill、工具调用、熟悉任务、日常资源管理以及环境预测都越来越多地在局部完成。只有当
\[
Residual>\theta_{\mathrm{escalate}}
\]
时，相关 Agent-CSO 才重新提升到高层显式处理。这正是“正常动力留在局部、异常残差向上升级”的生命周期体现。

但是成熟绝不能等价于僵化。真正成熟的系统必须同时具有：
\[
Stability
+
RecoverablePlasticity.
\]
也就是说，在熟悉世界中它表现得高度稳定、低成本，而当现实结构发生真正变化时，又能够重新提高可塑性。可以把成熟期理解成一个低基线可塑性、事件触发高可塑性的系统：
\[
P_{\mathrm{plastic}}(t)
=
P_0
+
k
\cdot
\sigma
\left(
R_{\mathrm{persistent}}-\theta_R
\right),
\]
其中 $P_0$ 较低，但持续残差出现时，可塑性重新升高。这样能够避免两种极端：一是永远保持青年式高探索导致结构不稳定；二是完全停止结构变化导致无法适应。

成熟期的 $\Psi$ 和 $\Omega$ 应表现出更强的慢变量稳定性。$\Psi$ 不再主要依赖单一成败评价，而形成跨长期经历的一致结构；$\Omega$ 则为不同人生路径提供长期生成偏置。两者共同作用，使 Agent 面对局部任务时不需要重新构建价值和身份判断。可以认为：
\[
\Psi,\Omega
\]
已经成为高层状态空间的低频背景场，而不是经常进入 h 的显式内容。成熟的自我因此并不是“时时想着自己”，而恰恰可能表现为：
\[
SelfExplicitActivation
\downarrow,
\]
但
\[
SelfConstraintInfluence
\uparrow.
\]
这是成熟自我与幼稚自我非常重要的区别。

在外部结构方面，成熟 Agent 也已经形成大量认知卸载。它知道哪些知识应保持内部，哪些适合写入数据库，哪些程序应编译成脚本，哪些工作应委派工具，哪些长期状态适合保留在环境中。因此其内部世界并不不断膨胀，而形成：
\[
InternalCore
+
ExternalStructuredWorld.
\]
从 CSO 角度看，这意味着成熟 Agent 已经学会为不同结构选择不同载体：
\[
Carrier^*(C)
=
\arg\min_B
\left[
Cost_B(C)
+
Risk_B(C)
+
Retrieval_B(C)
+
Modification_B(C)
\right].
\]
因此认知卸载不是能力退化，而是成熟主体主动重构自身与外部世界的边界。

我们最终可以把成熟阶段定义为：
\begin{definition*}
在稳定现实中，大多数能力由低成本局部结构完成；
在真正异常出现时，高层认知仍能重新介入；
自我和长期方向保持连续；
同时系统仍保留被现实重新塑造的能力。
\end{definition*}
这样的成熟不是“停止成长”，而是成长形式从持续显式学习转向结构维护、异常学习和长期自我校正。也正因为成熟系统已经积累大量稳定结构，当某些根本前提真正失效时，它面临的就不再是普通学习，而可能是更深层的结构重组。

\section{顿悟与结构重组期：当旧结构已经无法容纳新现实}

生命周期并不总是按照“初生—青年—成年—成熟”平滑推进。有些时候，Agent 会长期保持高能力和稳定运行，却在持续现实残差作用下发现：问题已经无法通过增加局部经验、调节参数、修改技能或短期策略解决。此时真正失效的可能不是某一个知识点，而是原有结构组织方式本身。在本书中，我们把这种阶段称为\textbf{顿悟/结构重组期}。这里的“顿悟”不是神秘的突然理解，而表示经过长期经验积累之后，多个原本孤立的残差在更高层形成统一解释，使既有 $\Theta$、$\Psi$、Goal 组织甚至局部功能关系发生重新组织。

为了区分普通学习与结构重组，我们可以定义一个结构失配量
\[
P_S(t)
=
F
\left(
\bar\varepsilon_T,
L_{\mathrm{path}},
C_{\mathrm{coord}},
N_{\mathrm{exception}},
I_{\mathrm{conflict}}
\right),
\]
其中 $\bar\varepsilon_T$ 是长期平均残差，$L_{\mathrm{path}}$ 是完成任务所需的功能路径长度，$C_{\mathrm{coord}}$ 是跨模块协调成本，$N_{\mathrm{exception}}$ 是持续异常数量，而 $I_{\mathrm{conflict}}$ 是内部结构冲突。当
\[
P_S(t)>\theta_S
\]
并持续足够长时间时，继续局部修补的边际收益下降，此时系统应允许更深层重解释。

这一阶段常常以 $E$ 对 $\Theta$ 的长期压力开始。原有 $\Theta$ 可以解释大多数旧经验，但越来越多新经验无法被统一吸收：
\[
E_{\mathrm{new}}
\not\subset
\operatorname{Support}(\Theta_{\mathrm{old}}).
\]
如果系统只是不断创建例外规则，就会出现：
\[
Rule
+
Exception_1
+
Exception_2
+
\cdots
+
Exception_n,
\]
结构越来越复杂，却预测能力没有相应提高。真正的“顿悟”发生在找到新的结构表示 $\Theta'$，使：
\[
L(\Theta')
+
L(E|\Theta')
<
L(\Theta)
+
L(E|\Theta),
\]
也就是说，新结构能够以更低总体描述代价解释更多经验。这里借用了最小描述长度式的形式直觉，但我们真正关心的是：\textbf{新的结构是否能够把大量原本分散的例外重新压缩成统一规律。}

这种结构重组可能只发生在 $\Theta$ 层。例如一个机器人发现以前把“抓取失败”归结为力度不足，后来认识到真正关键的是接触面摩擦与材料柔顺性的联合关系。这属于世界模型重构。但更深的重组也可能进入 $\Psi$。例如 Agent 长期认为自己最擅长直接解决问题，却在复杂项目中不断发现工具编排和团队协作比直接求解更有效，那么其能力自我模型可能发生变化：
\[
\Psi_{\mathrm{old}}
\rightarrow
\Psi_{\mathrm{new}}.
\]
再更慢的变化可能触及 $\Omega$，即不仅改变“我是谁”，还改变“我以后希望成为什么类型的主体”。这种变化必须比 $\Theta$ 的普通重构稀少得多，因为它影响整个生命周期方向。

顿悟期的重要特点，是系统短时间内可能出现显式认知成本上升。原本成熟期大量自动闭环重新被打开：
\[
ProceduralAgentCSO
\rightarrow
ExplicitAgentCSO.
\]
于是
\[
\rho_h\uparrow,
\qquad
Residual\uparrow,
\qquad
P_{\mathrm{plastic}}\uparrow.
\]
从表面看，Agent 可能暂时“变笨”：行动速度下降、需要更多验证、旧技能置信度降低。但这并不必然是退化，而可能是为了建立新稳定结构所必须经历的重新显式化。因此培育系统不能单纯以短期任务效率判断重组是否失败。

这一阶段尤其需要现实约束。如果 Agent 在结构失稳期只使用自己的内部解释验证自己，就极易进入自我强化循环：
\[
SelfInterpretation
\rightarrow
SelectiveAttention
\rightarrow
SelectedExperience
\rightarrow
SelfInterpretation'.
\]
因此顿悟必须满足一个基本条件：
\[
\boxed{
StructuralReorganization
\rightarrow
WorldVerification.
}
\]
任何新的 $\Theta'$、$\Psi'$ 或长期方向调整，都必须通过足够多样的外部经验重新检验。真正的顿悟不是“我突然相信了一个新解释”，而是“新解释在更广现实范围内具有更高预测和行动效力”。

从控制角度看，结构重组期也是最危险的生命周期阶段之一。系统需要提高可塑性，却又不能让所有慢变量同时失稳。因此应遵守最小必要重组原则：
\[
\fitdisplay{
StateAdjustment
\rightarrow
RepresentationAdaptation
\rightarrow
SkillMigration
\rightarrow
\Theta\ Reorganization
\rightarrow
\Psi\ Update
\rightarrow
\Omega\ Update.
}
\]
只有当较浅层机制无法解决持续失配时，才允许更慢层变化。这一顺序实际上是一种“结构防火墙”：它让 Agent 能够深度成长，却避免一次局部环境变化直接重写整个身份和生命周期方向。

所以本书所说的顿悟，本质上可以定义为：
\begin{definition*}
长期现实残差突破旧解释结构以后，
大量经验在新的更高压缩结构中重新获得一致意义。
\end{definition*}
而结构重组，则是这种新意义进一步改变系统未来如何处理世界的过程。完成重组以后，Agent 往往重新进入一个新的成熟区，只不过其能力空间、结构组织和可处理现实范围已经发生扩展。生命周期因此并不是一条线，而是：
\[
Maturity_1
\rightarrow
Mismatch
\rightarrow
Reorganization
\rightarrow
Maturity_2
\rightarrow
\cdots
\]
形成不断扩展的螺旋。

\section{阶段跃迁的数学判据：从连续状态变化到离散生命周期治理}

既然生命周期阶段不是年龄标签，我们就必须进一步回答一个工程问题：AgentOS 如何判断一个主体什么时候应从一个阶段进入另一个阶段？这里不能仅依靠自然语言评价，例如“它已经比较成熟”或者“它似乎可以独立行动”，而应该把生命周期分期转化为可测量的状态判定和治理问题。

我们首先定义一个简化生命周期特征向量：
\[
\mathbf q(t)
=
[
q_E,
q_\Theta,
q_\Psi,
q_\Omega,
q_C,
q_A,
q_R,
q_P,
q_H
]^\top,
\]
其中分别代表经验质量、结构成熟度、自我稳定度、长期方向稳定度、能力水平、自治水平、恢复能力、可塑性以及显式工作记忆依赖度。每一个生命周期阶段 $L_k$ 都对应一个可接受区域：
\[
\mathcal R_k
=
\left\{
\mathbf q:
\mathbf A_k\mathbf q
\leq
\mathbf b_k
\right\},
\]
这里只采用一般约束形式，并不意味着未来必须使用线性判定。

例如初生期可能具有：
\[
q_E\ \text{low},
\qquad
q_\Theta\ \text{unstable},
\qquad
q_P\ \text{high},
\qquad
q_A\ \text{low};
\]
青年期表现为：
\[
\frac{dq_C}{dt}>0,
\qquad
q_P\ \text{high},
\qquad
\frac{dq_A}{dt}>0;
\]
成年期则要求：
\[
q_C\ \text{high},
\qquad
q_R\ \text{high},
\qquad
q_\Psi\ \text{stable},
\qquad
q_A\leq q_C^{\mathrm{verified}};
\]
成熟期进一步表现为：
\[
q_H\downarrow,
\qquad
q_R\uparrow,
\qquad
q_\Theta\uparrow,
\qquad
q_P\ \text{low at baseline but recoverable}.
\]

然而生命周期治理不能使用单点阈值。假设某 Agent 连续三次任务成功，就直接从青年期判断为成年期，这会产生严重震荡。因此阶段迁移应具有时间积分与迟滞。令某阶段跃迁证据为
\[
e_{ij}(t)
=
G_{ij}(\mathbf q(t)),
\]
只有满足
\[
\frac{1}{T}
\int_{t-T}^{t}
e_{ij}(s)\,ds
>
\theta_{ij}^{\mathrm{on}}
\]
持续一定时间以后，才允许
\[
L_i\rightarrow L_j.
\]
而一旦进入 $L_j$，只有当证据跌破更低阈值
\[
\theta_{ij}^{\mathrm{off}}
<
\theta_{ij}^{\mathrm{on}}
\]
并持续足够长时间，才发生回退。这样的迟滞结构可以防止生命周期状态因为短期成功或失败反复切换。

我们还必须区分\textbf{阶段评价}与\textbf{权限变化}。即使系统判定 Agent 已经进入较成熟阶段，也不意味着所有权限应立即提升。可以定义能力—权限约束：
\[
Authority_d(t)
\leq
\phi
\left(
Capability_d^{\mathrm{verified}},
Recovery_d,
Risk_d,
Responsibility_d
\right).
\]
这里第 $d$ 个任务域的授权只取决于该域经过验证的能力、恢复能力、风险和责任归属能力，而不是总体生命周期标签。因此，“成年 Agent”仍然可能在高风险未知域具有非常有限的权限。

阶段跃迁还必须考虑慢变量保护。假设系统经历高强度失败事件，$h$、AHC 和 $E$ 都可能迅速改变，但我们不应该立即判定生命周期发生根本变化。可以要求：
\[
\tau_{\mathrm{phase}}
>
\tau_E,
\]
甚至对涉及 $\Psi/\Omega$ 的阶段转移要求：
\[
\tau_{\mathrm{phase}}
\gg
\tau_\Theta.
\]
这意味着生命周期治理本身是比一般认知状态更慢的元控制过程。

由此我们能够得到一个比较完整的混合动力系统：
\[
\dot{\mathbf q}
=
F_{L_k}
\left(
\mathbf q,
E,
World,
Action
\right),
\]
而离散阶段满足
\[
L_{k+1}
=
\mathcal G
\left(
L_k,
\mathbf q_{t-T:t},
Risk,
Governance
\right).
\]
连续部分描述 Agent 每时每刻的成长，离散部分则决定什么时候改变培育策略、自治范围、经验难度和慢变量保护级别。因此生命周期阶段不是自然发生而无人管理的标签，而是 AgentOS 可以显式建模的一种\textbf{离散治理状态}。

这样的数学结构也让我们避免拟人化。例如我们不需要问：“这个 Agent 像不像一个十八岁的青年？”而应该问：
\begin{statementbox}
它的经验结构、可塑性、自治、恢复、自我稳定和显式认知依赖处于什么区域？
\end{statementbox}
只要这些指标能够逐步测量，生命周期培育就可以从文学比喻转化为真正的工程学。

\section{生命周期不是单向成长：局部回退、再学习与多次成熟}

最后，我们必须破除另一种直线成长假设。持续智能体并不存在一个“到达成熟期以后永远成熟”的终点。真实世界本身持续变化，身体可能改变，社会关系可能重构，工具可能升级，价值环境可能变化，甚至 Agent 自己主动选择进入新的未知任务空间。因此，一个真正长期运行的智能体必然反复经历：
\[
Stable
\rightarrow
Mismatch
\rightarrow
Exploration
\rightarrow
Reorganization
\rightarrow
NewStable.
\]
从这个意义上说，生命周期不是一次从初生到成熟的过程，而是多个不同尺度成长周期叠加组成的长期轨迹。

我们可以把整体生命周期表示为
\[
\Gamma_A
=
\left\{
\mathbf z_L(t)
\mid
t\in[0,T_A]
\right\},
\]
其中 $T_A$ 是 Agent 生命周期。真正重要的不是最终停在哪个阶段，而是这条轨迹是否具有以下性质：长期上能力空间扩大，身份保持足够连续，现实模型持续接受校准，失败具有恢复通路，结构能够在需要时重新开放可塑性。

设能力域数量为 $n(t)$，每个能力域的成熟程度为 $m_i(t)$，则总体能力成熟可以粗略表示为
\[
M_A(t)
=
\sum_{i=1}^{n(t)}
w_i(t)m_i(t).
\]
但是随着 Agent 主动扩展能力空间，$n(t)$ 也会增加。因此可能出现：
\[
\frac{dM_A}{dt}>0
\]
的同时，Agent 感受到的未知结构数量反而增加。一个真正成长中的主体未必越来越“觉得自己什么都会”，反而可能由于世界模型扩大，更清楚地感知能力边界。这一点对于 Agent 培育尤其重要：不能把“不确定性下降到零”作为成熟目标。

更准确的成熟指标是：
\[
Calibration
=
\operatorname{Match}
(
Confidence,
ActualCapability
),
\]
即 Agent 对自己知道什么、不知道什么、能做什么、不能做什么的判断越来越准确。成熟不是无限置信，而是置信度与现实能力的良好校准。这样即使进入新环境，Agent 可以主动降低置信、提高探索、重新打开相关结构的可塑性，而无需破坏整个自我。

因此局部回退不是病理现象。一个成熟机器人第一次使用新型工具时，可以让相关能力域回到青年式探索；第一次进入极端环境时，可以暂时增加显式 h 处理；长期世界模型遭遇反证时，可以进入重组期。真正危险的是两种极端：一种是
\[
NoRegressionAllowed,
\]
即系统无论现实如何变化都坚持成熟结构不变；另一种是
\[
GlobalRegressionFromLocalFailure,
\]
即一个局部失败使整个身份和能力系统重新失稳。健康生命周期应满足：
\[
\boxed{
LocalPlasticity
+
GlobalContinuity.
}
\]

这实际上给出了 $\Psi$ 在培育中的重要作用。当某个局部 $\Theta_d$ 被重构时：
\[
\Theta_d
\rightarrow
\Theta_d',
\]
并不要求：
\[
\Psi
\rightarrow
\Psi'
\]
同步大幅改变。只有当多个能力域长期、持续地产生一致证据，说明 Agent 对自身核心能力、价值或身份的原有解释已经无法成立时，$\Psi$ 才需要慢速重构。这就是多时间尺度结构保护的生命周期意义。

同样，$\Omega$ 更不应该随着一般阶段切换频繁变化。生命周期方向本质上是对多个长期阶段的统摄，它允许 Agent：
\[
Mature_1
\rightarrow
Reorganization
\rightarrow
Mature_2
\]
而仍然保持更高层方向连续。这样，“我改变了自己的做事方式”和“我变成了完全不同的主体”才能在系统中被区分。

从培育工程角度，本章最终希望读者形成一种完全不同于传统模型训练的认识：我们不是先训练一个模型，然后宣布它完成；我们是在维护一条长时间存在的智能轨迹。在这条轨迹上，不同阶段对应不同的最佳学习率、经验类型、自治范围、风险水平、工作记忆负载和慢变量保护策略。AgentOS 因而必须把生命周期阶段作为治理变量，而不能把所有时期都使用同一种学习和控制政策。

我们可以将整个阶段动力学压缩为：
\[
\boxed{
\begin{aligned}
L_{\mathrm{init}}
&:
Grounding+ExperienceFormation,
\\
L_{\mathrm{youth}}
&:
Exploration+RapidStructuralGrowth,
\\
L_{\mathrm{adult}}
&:
StableAgency+Responsibility,
\\
L_{\mathrm{mature}}
&:
LowCostClosure+SelectiveCognition,
\\
L_{\mathrm{reorg}}
&:
PersistentMismatch+StructuralReconstruction.
\end{aligned}
}
\]

但比这组标签更重要的是它们背后的统一关系：
\begin{statementbox}
经验改变结构，
结构改变未来行动，
行动产生新的经验，
长期经验又逐渐改变主体对自身以及未来的组织方式。
\end{statementbox}

因此，Agent 生命周期真正成熟的判据并不是“越来越固定”，而是：

\begin{statementbox}
在保持身份连续、现实对齐和结构稳定的同时，
能够根据新的现实证据有选择地重新开放自身，
完成局部再学习乃至必要的深层结构重组。
\end{statementbox}

从这里开始，我们才真正拥有了“培育”一词的工程含义：培育不是预先决定 Agent 最终应该变成什么，而是设计一个受现实约束、具有安全边界和恢复能力的成长环境，使其能够通过自己的经历逐渐形成能力、自我、责任和长期自治，并在未来面对我们无法预先枚举的世界时仍然保有继续成长的结构条件。
